\section{DISCUSSION}

The experiments ask a more precise question than ``can a hybrid robot walk and
drive'': what behaviors can be obtained by reusing one actuated distal contact?
On the open-arc model, distal rotation alone traversed the $150\,$mm profile and
coordinated articulation was required by the tested controllers at $200\,$mm.
On flat ground, rotation alone was slower and more expensive than walking, while
the highest descriptive speed appeared when posture and rotation were commanded
together. These observations motivate contact reuse---support, quasi-rolling,
and candidate edge engagement from one link---but they do not yet attribute the
stair result uniquely to the opening. Equation~(\ref{eq:stairmoment}) explains
why: a wheel driven by axle torque does not share the singular force-only bound.
The matched open-arc/closed-wheel experiment is therefore the decisive test.

The costs are equally explicit. Coordinated arc locomotion produced
$6.23\times$ the slip and $3.18\times$ the mechanical work of walking, and
(\ref{eq:contactspeed}) explains the first-order mismatch: the executed
controller commands an activity-scaled surface speed rather than a contact-speed
feedback law. The commanded surface speed can therefore substantially exceed
body speed. The present model supports a speed/cost trade-off, not an
energy-efficiency claim. Rate control tied to measured body and contact velocity
is a controller improvement to test separately from the morphology ablation.

Four limitations bound how far these results generalize. First, the evidence is
a deterministic simulation: TPU compliance and friction, polymer infill mass,
gearbox efficiency, backlash, thermal limits and battery sag are approximated,
and the modeled mass remains a parameter until the assembled robot is weighed.
Second, the nominal conditions are single trials, so the reported differences
describe the three regimes on one model but carry no confidence intervals.
Third, the restrictions are not matched against each other: the walking and
rolling conditions use different cycle frequencies and stroke limits
(Sec.~\ref{sub:tripod}), so the comparison is between implemented regimes rather
than between equally parameterized controllers. Fourth, the edge-engagement
interpretation of the $150\,$mm result is inferred from geometry and from the
absence of articulation; it is not a direct morphology ablation. The locked
paired protocol of Sec.~\ref{sec:method} addresses the last three limitations in
simulation, but physical model identification remains necessary. The stair
supervisor is also given its geometry and performs no perception or foothold
search.

\section{CONCLUSION}

We presented an 18-actuator hexapod whose six open arc-shaped distal links are
reused for planted support, quasi-rolling contact and candidate stair-edge
engagement, together with a controller that commands tripod walking and
continuous arc rotation on the same joints at the same instant. In deterministic
MuJoCo trials the coordinated controller reached $0.200\,\mathrm{m\,s^{-1}}$
where articulated walking reached $0.117\,\mathrm{m\,s^{-1}}$, with substantially
greater work and slip; the open-arc stair restrictions traversed the $150\,$mm
profile without articulation and the $200\,$mm profile only with coordinated
articulation. Guarded transitions stayed upright and completed in under
$0.7\,$s. These single-run results define hypotheses and a reproducible operating
envelope, not causal morphology evidence. The accompanying matched-geometry,
paired-trial protocol makes the remaining open-arc claim falsifiable. Physical
speed, endurance, energy efficiency and stair reliability remain unvalidated.
